\documentclass[12pt,a4paper,oneside]{book}
\usepackage[utf8]{inputenc}
\usepackage[T1]{fontenc}
\usepackage{setspace}
\usepackage{indentfirst}
\usepackage{graphicx}
\usepackage{float}
\usepackage[table,xcdraw]{xcolor}
\usepackage{longtable}
\usepackage{booktabs}
\usepackage{hyperref}
\usepackage{listings}
\usepackage{xcolor}
\usepackage{makeidx}
\usepackage{amsmath}
\usepackage{amssymb}
\usepackage{geometry}
\usepackage{fancyhdr}
\usepackage{titlesec}
\usepackage{xurl}
\usepackage{enumitem}

% Page margins
\geometry{
    left=2.5cm,
    right=2.5cm,
    top=3cm,
    bottom=3cm,
    bindingoffset=0cm
}

% Hyperlink styling
\hypersetup{
    colorlinks=true,
    linkcolor=blue,
    filecolor=magenta,
    urlcolor=cyan,
    pdftitle={SIEWS+ 5.0 Technical Documentation},
    pdfpagemode=FullScreen,
}

% Code listing styling
\lstdefinestyle{codestyle}{
    backgroundcolor=\color{gray!10},
    basicstyle=\ttfamily\small,
    breaklines=true,
    frame=single,
    numbers=left,
    numbersep=5pt,
    tabsize=2,
}

\lstdefinestyle{pythonstyle}{
    language=Python,
    backgroundcolor=\color{gray!10},
    basicstyle=\ttfamily\small,
    breaklines=true,
    frame=single,
    numbers=left,
    numbersep=5pt,
}

\lstdefinestyle{jsonstyle}{
    language=JSON,
    backgroundcolor=\color{gray!10},
    basicstyle=\ttfamily\tiny,
    breaklines=true,
    frame=single,
    numbers=left,
    numbersep=5pt,
}

% Section formatting
\titleformat{\chapter}[display]
    {\normalfont\LARGE\bfseries}{\thechapter}{1em}{}
\titlespacing*{\chapter}{0pt}{20pt}{40pt}

% Header/Footer
\fancyhf{}
\fancyhead[LE,RO]{\thepage}
\fancyhead[RE]{\leftmark}
\fancyhead[LO]{\rightmark}
\pagestyle{fancy}

% Index
\makeindex

\begin{document}

% Title Page
\begin{titlepage}
    \centering
    \vspace*{3cm}

    {\Huge\bfseries SIEWS+ 5.0\par}
    \vspace{0.5cm}
    {\LARGE\bfseries Smart Integrated Early Warning System Plus\par}
    \vspace{1.5cm}

    {\LARGE Technical Documentation\par}
    \vspace{2cm}

    {\large Version 5.0\par}
    \vspace{1cm}

    {\large\textit{\textcopyright{} 2024-2026}\par}
    \vspace{2cm}

    \begin{minipage}{\linewidth}
        \centering
        AI-Powered Safety Monitoring System for Upstream Oil \& Gas Facilities
    \end{minipage}

    \vfill

\end{titlepage}

% Table of Contents
\pagenumbering{roman}
\tableofcontents
\newpage

\pagenumbering{arabic}
\onehalfspacing

% ============================================================================
% CHAPTER 1: PROJECT OVERVIEW
% ============================================================================
\chapter{Project Overview}

\section{Introduction}

\textbf{SIEWS+ 5.0 (Smart Integrated Early Warning System Plus)} is an advanced AI-powered safety monitoring system specifically designed for upstream oil and gas facilities, including offshore platforms, onshore processing areas, and industrial sites with high safety requirements.

The system provides real-time monitoring capabilities including:
\begin{itemize}
    \item Human presence detection using computer vision
    \item Zone-based restriction management via polygon overlays
    \item Personal Protective Equipment (PPE) violation detection
    \item Automated shutdown signal triggering for high-risk zone breaches
    \item Real-time WhatsApp notifications to supervisors
    \item Web-based monitoring dashboard
\end{itemize}

\section{Project Scope}

SIEWS+ 5.0 addresses critical safety monitoring requirements in hazardous industrial environments:

\begin{enumerate}
    \item \textbf{Real-time Person Detection}: Continuous monitoring of personnel presence using YOLOv8 deep learning models
    \item \textbf{PPE Compliance Monitoring}: Automated detection of safety equipment (helmets, vests, belts) usage
    \item \textbf{Zone Access Control}: Polygon-based restricted area management with dwell-time tracking
    \item \textbf{Environmental Hazard Detection}: Identification of road damage, open holes, and safety cones
    \item \textbf{Incident Documentation}: Automated snapshot capture and incident logging for compliance reporting
\end{enumerate}

\section{Technology Stack}

\begin{table}[H]
    \centering
    \caption{Technology Stack Overview}
    \begin{tabular}{|l|l|l|}
        \hline
        \textbf{Layer} & \textbf{Technology} & \textbf{Purpose} \\
        \hline
        Frontend & Next.js 14 (App Router) & Web Dashboard UI \\
        \hline
        Styling & Tailwind CSS 3.4 & Responsive Design \\
        \hline
        Backend & Python FastAPI & REST API, Streaming \\
        \hline
        Detection & YOLOv8n (Ultralytics) & Multi-stage Object Detection \\
        \hline
        Streaming & OpenCV + MJPEG & Video Frame Processing \\
        \hline
        Database & SQLite + SQLAlchemy & Persistent Storage \\
        \hline
        Realtime & FastAPI WebSocket & Live Alert Push \\
        \hline
        Notification & Fonnte API & WhatsApp Alerts \\
        \hline
        Face Rec. & dlib face\_recognition & Biometric Personnel ID \\
        \hline
        OCR & EasyOCR & Uniform Code Reading \\
        \hline
    \end{tabular}
\end{table}

\section{System Architecture Overview}

\begin{figure}[H]
    \centering
    \includegraphics[width=0.9\textwidth]{placeholder_arch}
    \caption{System Architecture Flow}
    \label{fig:system_architecture}
\end{figure}

The core data flow follows this sequence:
\begin{verbatim}
IP Camera → Backend (YOLOv8) → Zone Check (polygon) →
Alert + Shutdown → WhatsApp → Dashboard
\end{verbatim}

\section{Project Directory Structure}

\begin{verbatim}
migas-siews/
├── backend/
│   ├── main.py              # FastAPI entry point, all API endpoints
│   ├── stream.py            # StreamManager (detection pipeline)
│   ├── detector.py          # MultiStagePipeline (4-stage YOLO)
│   ├── polygon.py           # Ray-casting point-in-polygon
│   ├── zone_tracker.py     # PersonZoneTracker (dwell-time engine)
│   ├── violation_checker.py # ViolationChecker (PPE, zone, hazard)
│   ├── notifier.py          # WhatsApp via Fonnte API
│   ├── shutdown.py          # Shutdown signal handler
│   ├── models.py            # SQLAlchemy ORM models
│   ├── database.py          # SQLite engine, init_db()
│   ├── config.py            # Environment variables loader
│   ├── drawing.py           # OpenCV drawing utilities
│   ├── face_manager.py      # Face registration & recognition
│   ├── ocr_engine.py        # EasyOCR uniform code reader
│   ├── video_processor.py   # Async video upload processing
│   ├── detection_models.py # Class enums and thresholds
│   ├── requirements.txt    # Python dependencies
│   ├── .env                 # Environment configuration
│   └── static/
│       ├── snapshots/        # Violation screenshots
│       ├── faces/           # Registered face images
│       └── uploads/         # Uploaded video files
├── frontend/
│   ├── app/
│   │   ├── layout.tsx       # Root layout + Navbar
│   │   ├── page.tsx        # Redirects to /dashboard
│   │   ├── dashboard/      # Main monitoring page
│   │   ├── incidents/      # Incident log + filters
│   │   ├── zones/          # Zone management
│   │   ├── settings/      # System configuration
│   │   ├── analyze/       # Photo analysis page
│   │   ├── video-analysis/ # Video upload + processing
│   │   └── faces/         # Personnel biometric registry
│   ├── components/
│   │   ├── Navbar.tsx      # Top navigation bar
│   │   ├── VideoCanvas.tsx # MJPEG stream + polygon overlay
│   │   ├── ZoneEditor.tsx  # Zone list panel
│   │   ├── AlertFeed.tsx   # Live WebSocket alert sidebar
│   │   ├── AlertCard.tsx   # Individual alert card
│   │   ├── ShutdownBanner.tsx # Emergency shutdown banner
│   │   ├── ImageTester.tsx # Image upload component
│   │   └── LoadingScreen.tsx
│   ├── package.json
│   └── tailwind.config.ts
├── training/
│   ├── train_stage3_fire_smoke_v2.ipynb
│   ├── train_stage3_fire_smoke.py
│   └── requirements_train.txt
├── dataset/                 # Dataset storage
├── backend-go/              # Go-based backend (alternative)
├── camera-gateway/          # Camera gateway service
└── README.md
\end{verbatim}

% ============================================================================
% CHAPTER 2: BACKEND COMPONENTS
% ============================================================================
\chapter{Backend Components}

\section{Main Application (main.py)}

The \texttt{main.py} file serves as the FastAPI application entry point, handling all REST API endpoints, WebSocket connections, and CORS configuration.

\subsection{API Endpoints}

\begin{longtable}{|l|l|p{8cm}|}
    \hline
    \textbf{Method} & \textbf{Endpoint} & \textbf{Description} \\
    \hline
    GET & \texttt{/stream} & MJPEG video stream with YOLO overlay \\
    \hline
    GET & \texttt{/polygons} & List all zones \\
    \hline
    POST & \texttt{/polygons} & Create a new zone \\
    \hline
    PUT & \texttt{/polygons/\{id\}} & Update zone \\
    \hline
    DELETE & \texttt{/polygons/\{id\}} & Delete zone + cascade records \\
    \hline
    GET & \texttt{/alerts} & Paginated incident log with filters \\
    \hline
    POST & \texttt{/alerts/\{id\}/resolve} & Mark incident resolved \\
    \hline
    POST & \texttt{/alerts/\{id\}/false-positive} & Mark as false positive \\
    \hline
    GET & \texttt{/alerts/\{id\}/detections} & Per-object detection crops \\
    \hline
    POST & \texttt{/shutdown/trigger} & Manual shutdown signal \\
    \hline
    GET & \texttt{/settings} & Get all settings \\
    \hline
    POST & \texttt{/settings} & Update a setting \\
    \hline
    POST & \texttt{/settings/notify-test} & Send test WhatsApp \\
    \hline
    WS & \texttt{/ws/alerts} & Real-time alert push WebSocket \\
    \hline
    WS & \texttt{/ws/camera} & Browser camera → AI → annotated frame \\
    \hline
    GET & \texttt{/stats} & Dashboard statistics \\
    \hline
    GET & \texttt{/analytics/compliance} & PPE compliance statistics \\
    \hline
    POST & \texttt{/ai/analyze-image} & Analyze uploaded image \\
    \hline
    POST & \texttt{/analyze/image} & Upload image for multi-stage detection \\
    \hline
    POST & \texttt{/stream/simulate} & Inject static image into live stream \\
    \hline
    POST & \texttt{/stream/reset} & Reset to live camera \\
    \hline
    POST & \texttt{/video/upload} & Upload video for async processing \\
    \hline
    GET & \texttt{/video/jobs} & List video processing jobs \\
    \hline
    GET & \texttt{/video/jobs/\{id\}} & Get job status \\
    \hline
    GET & \texttt{/video/jobs/\{id\}/result} & Get full detection results \\
    \hline
    GET & \texttt{/video/annotated/\{id\}} & Serve annotated video \\
    \hline
    DELETE & \texttt{/video/jobs/\{id\}} & Delete a job \\
    \hline
    GET & \texttt{/faces} & List registered faces \\
    \hline
    POST & \texttt{/faces/register} & Register a new face \\
    \hline
    DELETE & \texttt{/faces/\{id\}} & Delete a face \\
    \hline
    POST & \texttt{/faces/train} & Re-load face database \\
    \hline
    POST & \texttt{/ai/generate-report} & LLM incident report generation \\
    \hline
    POST & \texttt{/ai/schedule} & Agentic scheduling (mock) \\
    \hline
    GET & \texttt{/health} & Health check \\
    \hline
\end{longtable}

\section{Database Models (models.py)}

The backend uses SQLAlchemy ORM for database management with SQLite as the database engine.

\subsection{Zone Model}

Represents restricted zone polygons on the monitoring area.

\begin{lstlisting}[style=pythonstyle]
class Zone(Base):
    __tablename__ = "zones"

    id = Column(Integer, primary_key=True, index=True)
    name = Column(String(100), nullable=False)
    vertices_json = Column(Text, nullable=False)  # JSON: [[x,y], ...]
    color = Column(String(7), default="#FF0000")
    active = Column(Boolean, default=True)
    risk_level = Column(String(10), default="high")  # "high" | "low"
    zone_type = Column(String(20), default="restricted")
    dwell_threshold_sec = Column(Integer, default=10)
    created_at = Column(DateTime, default=datetime.utcnow)
\end{lstlisting}

\subsection{Alert Model}

Incident log entry for detected violations.

\begin{lstlisting}[style=pythonstyle]
class Alert(Base):
    __tablename__ = "alerts"

    id = Column(Integer, primary_key=True, index=True)
    zone_id = Column(Integer, ForeignKey("zones.id"), nullable=True)
    confidence = Column(Float, default=0.0)
    snapshot_path = Column(String(255), nullable=True)
    timestamp = Column(DateTime, default=datetime.utcnow)
    shutdown_triggered = Column(Boolean, default=False)
    resolved = Column(Boolean, default=False)
    violation_type = Column(String(50), nullable=True)
    false_positive = Column(Boolean, default=False)
    ppe_detail = Column(Text, nullable=True)  # JSON
    person_name = Column(String(100), nullable=True)
    uniform_code = Column(String(50), nullable=True)
\end{lstlisting}

\subsection{DetectionLog Model}

Per-object crop for false positive review and audit.

\begin{lstlisting}[style=pythonstyle]
class DetectionLog(Base):
    __tablename__ = "detection_logs"

    id = Column(Integer, primary_key=True, index=True)
    alert_id = Column(Integer, ForeignKey("alerts.id"), nullable=True)
    class_name = Column(String(50), nullable=False)
    confidence = Column(Float, default=0.0)
    crop_path = Column(String(255), nullable=True)
    frame_number = Column(Integer, nullable=True)
    bbox_json = Column(Text, nullable=True)  # JSON: {x1,y1,x2,y2}
    is_false_positive = Column(Boolean, default=False)
    created_at = Column(DateTime, default=datetime.utcnow)
\end{lstlisting}

\subsection{Other Models}

\begin{itemize}
    \item \textbf{ShutdownLog}: Records shutdown events with trigger source (auto/manual)
    \item \textbf{Setting}: Key-value configuration store for system parameters
    \item \textbf{VideoJob}: Async video processing job tracking
    \item \textbf{ZoneOccupancy}: Zone dwell audit log for compliance reporting
\end{itemize}

\section{Core Detection Pipeline}

\subsection{StreamManager (stream.py)}

The StreamManager orchestrates the entire detection pipeline:

\begin{enumerate}
    \item Captures frames from IP camera at configured interval
    \item Runs multi-stage YOLO detection pipeline
    \item Checks for PPE violations using ViolationChecker
    \item Checks for zone violations using polygon algorithms
    \item Triggers alerts and shutdown signals when violations detected
    \item Generates MJPEG stream with overlay annotations
    \item Broadcasts alerts via WebSocket to connected dashboards
\end{enumerate}

\subsection{MultiStagePipeline (detector.py)}

Implements a 4-stage detection pipeline:

\begin{longtable}{|l|l|l|l|}
    \hline
    \textbf{Stage} & \textbf{Model} & \textbf{Classes} & \textbf{Purpose} \\
    \hline
    1 & yolov8n.pt & person & Base person detection \\
    \hline
    2 & best\_stage2\_labeled\_safety.pt & unknown, belt, helmet, vest & PPE compliance \\
    \hline
    3 & best\_stage3\_openhole.pt & barricade, hard\_hat, safety\_cone, open\_hole, vest\_env & Environmental hazards \\
    \hline
    4 & best\_jalan\_berlubang.pt & lubang, retak, tambalan & Road damage detection \\
    \hline
\end{longtable}

\section{Polygon Algorithms (polygon.py)}

\subsection{Ray-Casting Point-in-Polygon}

The system uses the ray-casting algorithm to determine if a point is inside a polygon:

\begin{lstlisting}[style=pythonstyle]
def point_in_polygon(point: Tuple[float, float], polygon: List[List[float]]) -> bool:
    """Ray-casting algorithm to determine if a point is inside a polygon.
    All coordinates are normalized (0.0 - 1.0)."""
    x, y = point
    n = len(polygon)
    inside = False

    j = n - 1
    for i in range(n):
        xi, yi = polygon[i]
        xj, yj = polygon[j]

        if ((yi > y) != (yj > y)) and (x < (xj - xi) * (y - yi) / (yj - yi) + xi):
            inside = not inside
        j = i

    return inside
\end{lstlisting}

All polygon coordinates are normalized (0.0 to 1.0) for resolution independence.

\section{Violation Checking (violation\_checker.py)}

The ViolationChecker validates detected objects against safety rules:

\subsection{PPE Violations}
\begin{itemize}
    \item NO\_HELMET: Person detected without helmet
    \item NO\_VEST: Person detected without safety vest
    \item NO\_BELT: Person detected without safety belt
\end{itemize}

\subsection{Zone Violations}
\begin{itemize}
    \item Checks person (head, chest, feet points) inside restricted zones
    \item Triggers shutdown for high-risk zones
    \item Logs warning for low-risk zones
\end{itemize}

\subsection{Hazard Violations}
\begin{itemize}
    \item Person center overlapping environmental hazard (open hole, road damage)
    \item Generates critical alert with snapshot
\end{itemize}

\section{Person Tracking (zone\_tracker.py)}

The PersonZoneTracker maintains stable person identities across frames:

\begin{itemize}
    \item Uses centroid proximity matching for track association
    \item Generates unique 8-character hex UUIDs for track IDs
    \item Tracks dwell time in zones with configurable thresholds
    \item Emits events: zone\_entry, dwell\_warning, dwell\_critical, zone\_exit
\end{itemize}

\section{Notification System (notifier.py)}

WhatsApp notifications via Fonnte API:

\begin{lstlisting}[style=pythonstyle]
def send_whatsapp(message: str, recipient: str) -> dict:
    """Send single WhatsApp message via Fonnte API."""
    url = "https://fonnte.com/api/send_message.php"
    data = {
        "token": FONNTE_TOKEN,
        "target": recipient,
        "message": message,
    }
    response = requests.post(url, data=data)
    return response.json()
\end{lstlisting}

\section{Face Recognition (face\_manager.py)}

\begin{itemize}
    \item Uses dlib face\_recognition (128-dim HOG encoding)
    \item OpenCV Haar cascade fallback when dlib unavailable
    \item Registers faces with name, code, and phone number
    \item Stores encodings in face\_db.json for persistence
    \item Singleton instance: face\_manager
\end{itemize}

\section{OCR Engine (ocr\_engine.py)}

\begin{itemize}
    \item Uses EasyOCR for uniform code reading from person chest area
    \item Applies CLAHE, denoising, and sharpening preprocessing
    \item Fuzzy-matches against registered codes from face database
    \item Graceful fallback when EasyOCR unavailable
\end{itemize}

% ============================================================================
% CHAPTER 3: FRONTEND COMPONENTS
% ============================================================================
\chapter{Frontend Components}

\section{Next.js Application Structure}

The frontend is built with Next.js 14 App Router with Tailwind CSS for styling.

\subsection{Page Routes}

\begin{longtable}{|l|l|p{8cm}|}
    \hline
    \textbf{Page} & \textbf{Path} & \textbf{Purpose} \\
    \hline
    Root & \texttt{/} & Redirects to \texttt{/dashboard} \\
    \hline
    Dashboard & \texttt{/dashboard} & Main monitoring - live stream, zone editor, alerts \\
    \hline
    Zones & \texttt{/zones} & Zone management without canvas drawing \\
    \hline
    Incidents & \texttt{/incidents} & Historical alert log with CSV export \\
    \hline
    Settings & \texttt{/settings} & Camera, inference, WhatsApp configuration \\
    \hline
    Analyze Photo & \texttt{/analyze} & Single image upload + detection analysis \\
    \hline
    Video Analysis & \texttt{/video-analysis} & Video upload, job list, per-frame results \\
    \hline
    Faces & \texttt{/faces} & Personnel biometric registry \\
    \hline
\end{longtable}

\section{Key Components}

\subsection{VideoCanvas.tsx}

The primary monitoring component that displays the MJPEG stream with polygon overlay capabilities:

\textbf{Features:}
\begin{itemize}
    \item MJPEG stream display from \texttt{/stream} endpoint
    \item Canvas overlay for zone polygon drawing
    \item Click-to-place vertices, double-click to close polygon
    \item Alert flash animation (red border + overlay text)
    \item Camera offline overlay with reconnection logic
    \item LIVE indicator with ping animation
\end{itemize}

\textbf{VideoCanvas Interface:}
\begin{lstlisting}[style=pythonstyle]
interface ZoneData {
  id: number;
  name: string;
  vertices: number[][];  // Normalized [x, y] coordinates
  color: string;
  active: boolean;
  risk_level: string;
}
\end{lstlisting}

\subsection{ZoneEditor.tsx}

Zone management panel with the following features:
\begin{itemize}
    \item Lists all zones with active/inactive toggle
    \item "Tambah Zona" button to start drawing mode
    \item Delete zone button
    \item Risk level badge display
    \item Vertex count display
\end{itemize}

\subsection{AlertFeed.tsx}

Real-time alert sidebar component:
\begin{itemize}
    \item WebSocket connection to \texttt{/ws/alerts}
    \item Real-time alert list (max 30, newest first)
    \item Audio beep on new alerts (oscillator-based, different tone for high/low risk)
    \item Resolve and False Positive actions
    \item Connection status indicator
\end{itemize}

\subsection{Navbar.tsx}

Fixed navigation bar with:
\begin{itemize}
    \item Logo and navigation links
    \item Active link highlighting (amber color)
    \item Real-time stats chips: camera status, zone count, today's alerts
    \item Unresolved alert badge on "Incidents" link
\end{itemize}

\subsection{ShutdownBanner.tsx}

Emergency shutdown alert display:
\begin{itemize}
    \item Full-width emergency shutdown banner
    \item Displays zone name and timestamp (WIB timezone)
    \item "Confirm \& Reset" button to resolve alert
\end{itemize}

\section{State Management}

The frontend uses React's built-in state management with hooks:
\begin{itemize}
    \item \texttt{useState} for local component state
    \item \texttt{useEffect} for side effects (WebSocket, polling)
    \item \texttt{useCallback} for memoized callbacks
    \item API calls to FastAPI backend via \texttt{NEXT\_PUBLIC\_API\_URL}
\end{itemize}

\section{Styling}

All styling is done with Tailwind CSS utility classes. Custom theme colors:
\begin{itemize}
    \item Primary: Amber (\texttt{amber-500})
    \item Danger: Red (\texttt{red-600})
    \item Success: Green (\texttt{green-500})
    \item Info: Blue (\texttt{blue-500})
\end{itemize}

% ============================================================================
% CHAPTER 4: DATABASE SCHEMA
% ============================================================================
\chapter{Database Schema}

\section{Entity Relationship Diagram}

The database consists of 7 primary tables with the following relationships:

\begin{itemize}
    \item \texttt{alerts} references \texttt{zones} (zone\_id)
    \item \texttt{detection\_logs} references \texttt{alerts} (alert\_id)
    \item \texttt{shutdown\_logs} references \texttt{zones} (zone\_id)
    \item \texttt{video\_jobs} is standalone
    \item \texttt{zone\_occupancy} references \texttt{zones} (zone\_id)
    \item \texttt{settings} is a key-value store
\end{itemize}

\section{Table Definitions}

\subsection{zones}

\begin{tabular}{|l|l|l|p{6cm}|}
    \hline
    \textbf{Column} & \textbf{Type} & \textbf{Default} & \textbf{Description} \\
    \hline
    id & INTEGER & PK & Primary key \\
    \hline
    name & VARCHAR(100) & NOT NULL & Zone display name \\
    \hline
    vertices\_json & TEXT & NOT NULL & JSON list of normalized [x,y] pairs \\
    \hline
    color & VARCHAR(7) & \#FF0000 & Hex color code \\
    \hline
    active & BOOLEAN & TRUE & Zone enabled/disabled \\
    \hline
    risk\_level & VARCHAR(10) & high & "high" or "low" \\
    \hline
    zone\_type & VARCHAR(20) & restricted & Zone classification \\
    \hline
    dwell\_threshold\_sec & INTEGER & 10 & Dwell time before alert \\
    \hline
    created\_at & DATETIME & now() & Creation timestamp \\
    \hline
\end{tabular}

\subsection{alerts}

\begin{tabular}{|l|l|l|p{6cm}|}
    \hline
    \textbf{Column} & \textbf{Type} & \textbf{Default} & \textbf{Description} \\
    \hline
    id & INTEGER & PK & Primary key \\
    \hline
    zone\_id & INTEGER & FK & Reference to zones.id \\
    \hline
    confidence & FLOAT & 0.0 & Detection confidence \\
    \hline
    snapshot\_path & VARCHAR(255) & NULL & Path to violation screenshot \\
    \hline
    timestamp & DATETIME & now() & Alert timestamp \\
    \hline
    shutdown\_triggered & BOOLEAN & FALSE & Shutdown signal sent \\
    \hline
    resolved & BOOLEAN & FALSE & Alert resolved status \\
    \hline
    violation\_type & VARCHAR(50) & NULL & Type of violation detected \\
    \hline
    false\_positive & BOOLEAN & FALSE & Marked as false positive \\
    \hline
    ppe\_detail & TEXT & NULL & JSON with PPE details \\
    \hline
    person\_name & VARCHAR(100) & NULL & Recognized person name \\
    \hline
    uniform\_code & VARCHAR(50) & NULL & OCR-detected uniform code \\
    \hline
\end{tabular}

\subsection{detection\_logs}

\begin{tabular}{|l|l|l|p{6cm}|}
    \hline
    \textbf{Column} & \textbf{Type} & \textbf{Default} & \textbf{Description} \\
    \hline
    id & INTEGER & PK & Primary key \\
    \hline
    alert\_id & INTEGER & FK & Reference to alerts.id \\
    \hline
    class\_name & VARCHAR(50) & NOT NULL & Detected class name \\
    \hline
    confidence & FLOAT & 0.0 & Detection confidence \\
    \hline
    crop\_path & VARCHAR(255) & NULL & Path to cropped image \\
    \hline
    frame\_number & INTEGER & NULL & Source video frame number \\
    \hline
    bbox\_json & TEXT & NULL & JSON bounding box coords \\
    \hline
    is\_false\_positive & BOOLEAN & FALSE & Marked as false positive \\
    \hline
    created\_at & DATETIME & now() & Creation timestamp \\
    \hline
\end{tabular}

\subsection{shutdown\_logs}

\begin{tabular}{|l|l|l|p{6cm}|}
    \hline
    \textbf{Column} & \textbf{Type} & \textbf{Default} & \textbf{Description} \\
    \hline
    id & INTEGER & PK & Primary key \\
    \hline
    zone\_id & INTEGER & FK & Reference to zones.id \\
    \hline
    trigger\_source & VARCHAR(20) & NOT NULL & "auto" or "manual" \\
    \hline
    triggered\_at & DATETIME & now() & Shutdown timestamp \\
    \hline
\end{tabular}

\subsection{settings}

\begin{tabular}{|l|l|l|p{6cm}|}
    \hline
    \textbf{Column} & \textbf{Type} & \textbf{Default} & \textbf{Description} \\
    \hline
    id & INTEGER & PK & Primary key \\
    \hline
    key & VARCHAR(100) & UNIQUE & Setting key name \\
    \hline
    value & TEXT & NULL & Setting value \\
    \hline
\end{tabular}

\subsection{video\_jobs}

\begin{tabular}{|l|l|l|p{6cm}|}
    \hline
    \textbf{Column} & \textbf{Type} & \textbf{Default} & \textbf{Description} \\
    \hline
    id & INTEGER & PK & Primary key \\
    \hline
    filename & VARCHAR(255) & NOT NULL & Original filename \\
    \hline
    file\_path & VARCHAR(255) & NOT NULL & Stored file path \\
    \hline
    status & VARCHAR(20) & pending & pending/processing/done/failed \\
    \hline
    progress & INTEGER & 0 & Progress percentage 0-100 \\
    \hline
    total\_frames & INTEGER & 0 & Total frames in video \\
    \hline
    processed\_frames & INTEGER & 0 & Frames processed so far \\
    \hline
    result\_json & TEXT & NULL & Full detection results JSON \\
    \hline
    annotated\_video\_path & VARCHAR(255) & NULL & Path to output video \\
    \hline
    error\_message & TEXT & NULL & Error description if failed \\
    \hline
    created\_at & DATETIME & now() & Job creation time \\
    \hline
    completed\_at & DATETIME & NULL & Job completion time \\
    \hline
\end{tabular}

\subsection{zone\_occupancy}

\begin{tabular}{|l|l|l|p{6cm}|}
    \hline
    \textbf{Column} & \textbf{Type} & \textbf{Default} & \textbf{Description} \\
    \hline
    id & INTEGER & PK & Primary key \\
    \hline
    zone\_id & INTEGER & FK & Reference to zones.id \\
    \hline
    track\_id & VARCHAR(20) & NOT NULL & 8-char hex UUID for person \\
    \hline
    entry\_time & DATETIME & NOT NULL & Zone entry timestamp \\
    \hline
    exit\_time & DATETIME & NULL & Zone exit timestamp \\
    \hline
    dwell\_sec & INTEGER & NULL & Total dwell duration \\
    \hline
    alert\_level & VARCHAR(20) & none & none/warning/critical \\
    \hline
\end{tabular}

% ============================================================================
% CHAPTER 5: DETECTION MODELS
% ============================================================================
\chapter{Detection Models and Classes}

\section{Detection Class Definitions (detection\_models.py)}

\subsection{PPE Classes (Stage 2)}

\begin{tabular}{|l|l|l|}
    \hline
    \textbf{Index} & \textbf{Name} & \textbf{Description} \\
    \hline
    0 & UNKNOWN & Unknown object class \\
    \hline
    1 & BELT & Safety belt/buckle \\
    \hline
    2 & HELMET & Safety helmet/hard hat \\
    \hline
    3 & VEST & Safety vest \\
    \hline
\end{tabular}

\subsection{Environment Classes (Stage 3)}

\begin{tabular}{|l|l|l|}
    \hline
    \textbf{Index} & \textbf{Name} & \textbf{Description} \\
    \hline
    0 & BARRICADE & Safety barricade \\
    \hline
    1 & HARD\_HAT & Worker wearing hard hat \\
    \hline
    2 & SAFETY\_CONE & Safety cone marker \\
    \hline
    3 & OPEN\_HOLE & Open hole/excavation \\
    \hline
    4 & VEST\_ENV & Worker wearing safety vest (env) \\
    \hline
\end{tabular}

\subsection{Road Damage Classes (Stage 4)}

\begin{tabular}{|l|l|l|}
    \hline
    \textbf{Index} & \textbf{Name} & \textbf{Description} \\
    \hline
    0 & LUBANG & Pothole \\
    \hline
    1 & RETAK & Crack \\
    \hline
    2 & TAMBALAN & Patch/filled area \\
    \hline
\end{tabular}

\section{Confidence Thresholds}

\begin{tabular}{|l|l|l|}
    \hline
    \textbf{Detection Type} & \textbf{Threshold} & \textbf{Description} \\
    \hline
    PERSON\_CONFIDENCE\_THRESHOLD & 0.50 & Minimum person detection confidence \\
    \hline
    PPE\_CONFIDENCE\_THRESHOLD & 0.35 & Minimum PPE detection confidence \\
    \hline
    ENV\_CONFIDENCE\_THRESHOLD & 0.40 & Minimum environment hazard confidence \\
    \hline
    ROAD\_CONFIDENCE\_THRESHOLD & 0.40 & Minimum road damage confidence \\
    \hline
    SAFETY\_CONE\_CONFIDENCE & 0.50 & Minimum safety cone confidence \\
    \hline
    PPE\_IOU\_THRESHOLD & 0.10 & IoU threshold for PPE-person assignment \\
    \hline
\end{tabular}

\section{Model Files}

\begin{longtable}{|l|l|l|l|}
    \hline
    \textbf{Model File} & \textbf{Stage} & \textbf{Classes} & \textbf{Use Case} \\
    \hline
    best\_stage2\_labeled\_safety.pt & 2 (PPE) & unknown, belt, helmet, vest & PPE compliance \\
    \hline
    best\_stage3\_openhole.pt & 3 (Env) & barricade, hard\_hat, safety\_cone, open\_hole, vest\_env & Environmental hazards \\
    \hline
    best\_jalan\_berlubang.pt & 4 (Road) & lubang, retak, tambalan & Road damage detection \\
    \hline
    stage2\_ppe\_harness.pt & 2 (PPE) & harness-specific & PPE variant \\
    \hline
    stage3\_environment.pt & 3 (Env) & environment-specific & Environment variant \\
    \hline
    stage4\_infrastructure.pt & 4 (Infra) & infrastructure-specific & Infrastructure variant \\
    \hline
\end{longtable}

\section{Training Configuration}

Training uses YOLOv8n as the base model with the following hyperparameters:

\begin{longtable}{|l|l|l|p{6cm}|}
    \hline
    \textbf{Parameter} & \textbf{Value} & \textbf{Type} & \textbf{Description} \\
    \hline
    epochs & 80 & int & Training epochs \\
    \hline
    batch & 8-16 & int & Batch size (auto-adjusted for GPU VRAM) \\
    \hline
    imgsz & 640 & int & Input image size \\
    \hline
    patience & 15 & int & Early stopping patience \\
    \hline
    freeze & 10 & int & Frozen layers (transfer learning) \\
    \hline
    optimizer & AdamW & string & Optimizer type \\
    \hline
    lr0 & 0.001 & float & Initial learning rate \\
    \hline
    lrf & 0.01 & float & Final learning rate (cosine decay) \\
    \hline
    hsv\_h & 0.03 & float & Hue augmentation \\
    \hline
    hsv\_s & 0.9 & float & Saturation augmentation \\
    \hline
    hsv\_v & 0.5 & float & Value/brightness augmentation \\
    \hline
    mosaic & 1.0 & float & 4-image mosaic augmentation \\
    \hline
    mixup & 0.1 & float & Mixup augmentation probability \\
    \hline
    flipud & 0.1 & float & Vertical flip probability \\
    \hline
    degrees & 5.0 & float & Rotation augmentation \\
    \hline
    scale & 0.6 & float & Scale augmentation \\
    \hline
    erasing & 0.3 & float & Random erasing probability \\
    \hline
\end{longtable}

% ============================================================================
% CHAPTER 6: CONFIGURATION
% ============================================================================
\chapter{Configuration}

\section{Backend Environment (.env)}

\begin{lstlisting}[style=jsonstyle]
CAMERA_SOURCE=0
FACILITY_NAME=Offshore Platform A
CONFIDENCE_THRESHOLD=0.3
DETECTION_INTERVAL=3
NOTIFY_COOLDOWN=300
FONNTE_TOKEN=
DEFAULT_RECIPIENTS=
\end{lstlisting}

\section{Backend Requirements (requirements.txt)}

\begin{lstlisting}[style=pythonstyle]
fastapi==0.115.0
uvicorn==0.30.0
opencv-python==4.10.0.84
ultralytics==8.4.37
sqlalchemy==2.0.35
aiofiles==24.1.0
python-dotenv==1.0.1
httpx==0.27.2
python-multipart==0.0.9
websockets==12.0
face-recognition==1.3.0
easyocr==1.7.0
dlib==19.24.0
numpy>=1.24.0
pillow>=10.0.0
\end{lstlisting}

\section{Frontend Configuration (package.json)}

\begin{lstlisting}[style=jsonstyle]
{
  "name": "siews-frontend",
  "version": "5.0.0",
  "private": true,
  "scripts": {
    "dev": "next dev",
    "build": "next build",
    "start": "next start",
    "lint": "next lint"
  },
  "dependencies": {
    "next": "14.2.35",
    "react": "^18",
    "react-dom": "^18"
  },
  "devDependencies": {
    "@types/node": "^20",
    "@types/react": "^18",
    "tailwindcss": "^3.4.1",
    "typescript": "^5"
  }
}
\end{lstlisting}

\section{Tailwind Configuration (tailwind.config.ts)}

\begin{lstlisting}[style=pythonstyle]
import type { Config } from "tailwindcss";

const config: Config = {
  content: [
    "./pages/**/*.{js,ts,jsx,tsx,mdx}",
    "./components/**/*.{js,ts,jsx,tsx,mdx}",
    "./app/**/*.{js,ts,jsx,tsx,mdx}",
  ],
  theme: {
    extend: {
      colors: {
        primary: {
          50: "#fffbeb",
          100: "#fef3c7",
          200: "#fde68a",
          300: "#fcd34d",
          400: "#fbbf24",
          500: "#f59e0b",
          600: "#d97706",
          700: "#b45309",
          800: "#92400e",
          900: "#78350f",
        },
      },
    },
  },
  plugins: [],
};
export default config;
\end{lstlisting}

% ============================================================================
% CHAPTER 7: INSTALLATION AND SETUP
% ============================================================================
\chapter{Installation and Setup}

\section{System Requirements}

\begin{itemize}
    \item Python 3.10+
    \item Node.js 18+
    \item 8GB+ RAM (16GB recommended for training)
    \item NVIDIA GPU with 6GB+ VRAM (for detection)
    \item Internet connection for API services
\end{itemize}

\section{Backend Setup}

\subsection{Clone and Install}

\begin{verbatim}
# Clone repository
git clone https://github.com/username/migas-siews.git
cd migas-siews/backend

# Create virtual environment
python -m venv venv
source venv/bin/activate  # Linux/Mac
venv\Scripts\activate     # Windows

# Install dependencies
pip install -r requirements.txt

# Configure environment
cp .env.example .env
# Edit .env with your configuration
\end{verbatim}

\subsection{Run Backend}

\begin{verbatim}
# Development
uvicorn main:app --reload --host 0.0.0.0 --port 8001

# Production
uvicorn main:app --host 0.0.0.0 --port 8001 --workers 4
\end{verbatim}

\section{Frontend Setup}

\subsection{Install and Run}

\begin{verbatim}
# Navigate to frontend
cd ../frontend

# Install dependencies
npm install

# Run development server
npm run dev

# Build for production
npm run build
npm start
\end{verbatim}

\section{Dataset Preparation}

\subsection{Dataset Structure}

\begin{verbatim}
dataset/
├── Fire and Smoke.v14i.yolo26.zip
├── jalan berlubang.v3-tambahan-dataset.yolo26/
├── stage3/
│   ├── train/
│   │   ├── images/
│   │   └── labels/
│   ├── val/
│   │   ├── images/
│   │   └── labels/
│   └── test/
│       ├── images/
│       └── labels/
└── stage3_environment.yaml
\end{verbatim}

\subsection{YAML Configuration}

\begin{lstlisting}[style=pythonstyle]
# stage3_fire_smoke.yaml
path: /path/to/dataset/stage3
train: train/images
val: val/images
test: test/images

nc: 2
names: ['fire', 'smoke']
\end{lstlisting}

\section{Training}

\subsection{Local Training}

\begin{verbatim}
cd training
python train_stage3_fire_smoke.py
\end{verbatim}

\subsection{Kaggle Training}

\begin{enumerate}
    \item Upload training notebook to Kaggle
    \item Attach dataset from \texttt{/kaggle/input}
    \item Enable GPU P100/T4
    \item Run All cells
    \item Download \texttt{best\_stage3\_fire\_smoke.pt} from Output
\end{enumerate}

\section{Model Deployment}

\begin{verbatim}
# Copy trained model to backend/models/
cp best_stage3_fire_smoke.pt ../backend/models/

# Restart backend to load new model
uvicorn main:app --reload
\end{verbatim}

% ============================================================================
% CHAPTER 8: API REFERENCE
% ============================================================================
\chapter{API Reference}

\section{REST Endpoints}

\subsection{Zone Management}

\subsubsection{GET /polygons}

Retrieve all zones.

\textbf{Response:}
\begin{lstlisting}[style=jsonstyle]
{
  "zones": [
    {
      "id": 1,
      "name": "Restricted Area A",
      "vertices": [[0.1, 0.2], [0.3, 0.4], ...],
      "color": "#FF0000",
      "active": true,
      "risk_level": "high",
      "zone_type": "restricted",
      "dwell_threshold_sec": 10
    }
  ]
}
\end{lstlisting}

\subsubsection{POST /polygons}

Create a new zone.

\textbf{Request Body:}
\begin{lstlisting}[style=jsonstyle]
{
  "name": "Restricted Area A",
  "vertices": [[0.1, 0.2], [0.3, 0.4], [0.5, 0.6]],
  "color": "#FF0000",
  "risk_level": "high",
  "zone_type": "restricted",
  "dwell_threshold_sec": 10
}
\end{lstlisting}

\subsection{Alert Management}

\subsubsection{GET /alerts}

Retrieve paginated alerts.

\textbf{Query Parameters:}
\begin{itemize}
    \item \texttt{page}: Page number (default: 1)
    \item \texttt{limit}: Items per page (default: 20)
    \item \texttt{resolved}: Filter by resolved status
    \item \texttt{violation\_type}: Filter by violation type
    \item \texttt{zone\_id}: Filter by zone ID
    \item \texttt{start\_date}: Filter by start date
    \item \texttt{end\_date}: Filter by end date
\end{itemize}

\textbf{Response:}
\begin{lstlisting}[style=jsonstyle]
{
  "alerts": [...],
  "total": 100,
  "page": 1,
  "limit": 20
}
\end{lstlisting}

\section{WebSocket Endpoints}

\subsection{/ws/alerts}

Real-time alert push to connected clients.

\textbf{Server Message:}
\begin{lstlisting}[style=jsonstyle]
{
  "type": "alert",
  "data": {
    "id": 123,
    "zone_name": "Restricted Area A",
    "violation_type": "ZONE_VIOLATION",
    "confidence": 0.85,
    "timestamp": "2024-01-15T10:30:00Z",
    "snapshot_path": "/static/snapshots/alert_123.jpg",
    "shutdown_triggered": true
  }
}
\end{lstlisting}

\subsection{/ws/camera}

Browser camera → AI → annotated frame pipeline.

\textbf{Client sends:} Base64 encoded frame
\textbf{Server receives:} Annotated frame with bounding boxes

\section{Detection Classes Reference}

\subsection{PPE Detection Classes}

\begin{tabular}{|l|l|l|}
    \hline
    \textbf{Class ID} & \textbf{Name} & \textbf{Associated Violation} \\
    \hline
    0 & UNKNOWN & None \\
    \hline
    1 & BELT & NO\_BELT \\
    \hline
    2 & HELMET & NO\_HELMET \\
    \hline
    3 & VEST & NO\_VEST \\
    \hline
\end{tabular}

\subsection{Violation Types}

\begin{tabular}{|l|l|l|}
    \hline
    \textbf{Violation Type} & \textbf{Severity} & \textbf{Action} \\
    \hline
    NO\_HELMET & High & Alert + Log \\
    \hline
    NO\_VEST & High & Alert + Log \\
    \hline
    NO\_BELT & Medium & Alert + Log \\
    \hline
    ZONE\_VIOLATION & High & Alert + Shutdown + WhatsApp \\
    \hline
    DWELL\_WARNING & Medium & Alert + Log \\
    \hline
    DWELL\_CRITICAL & High & Alert + Shutdown + WhatsApp \\
    \hline
    HAZARD\_PROXIMITY & High & Alert + Log \\
    \hline
\end{tabular}

% ============================================================================
% CHAPTER 9: TROUBLESHOOTING
% ============================================================================
\chapter{Troubleshooting}

\section{Common Issues}

\subsection{Camera Not Displaying}

\textbf{Symptoms:} VideoCanvas shows "Camera Offline" overlay.

\textbf{Solutions:}
\begin{enumerate}
    \item Check camera source configuration in \texttt{.env}
    \item Verify IP camera is accessible and streaming
    \item Check if MJPEG stream endpoint responds: \texttt{GET /stream}
    \item Restart backend service
\end{enumerate}

\subsection{Detection Not Working}

\textbf{Symptoms:} No bounding boxes on stream, no person detection.

\textbf{Solutions:}
\begin{enumerate}
    \item Verify YOLO model files exist in \texttt{backend/models/}
    \item Check model filenames match configuration in \texttt{detector.py}
    \item Ensure GPU is available and CUDA is installed
    \item Check confidence thresholds are appropriate
    \item Review backend logs for detection errors
\end{enumerate}

\subsection{WhatsApp Notifications Not Sending}

\textbf{Symptoms:} Alerts logged but no WhatsApp message received.

\textbf{Solutions:}
\begin{enumerate}
    \item Verify FONNTE\_TOKEN is set correctly in \texttt{.env}
    \item Check DEFAULT\_RECIPIENTS contains valid phone numbers
    \item Test notification manually via \texttt{POST /settings/notify-test}
    \item Check Fonnte API status and quota
\end{enumerate}

\subsection{Zone Polygon Drawing Not Working}

\textbf{Symptoms:} Clicking on canvas does not place vertices.

\textbf{Solutions:}
\begin{enumerate}
    \item Ensure "Tambah Zona" drawing mode is activated
    \item Check browser console for JavaScript errors
    \item Verify WebSocket connection is established
    \item Clear browser cache and reload
\end{enumerate}

\subsection{Training Fails}

\textbf{Symptoms:} Training notebook errors or poor mAP results.

\textbf{Solutions:}
\begin{enumerate}
    \item Verify dataset structure is correct (train/val/images + labels)
    \item Check label files for correct class IDs (0=fire, 1=smoke)
    \item Ensure YAML \texttt{names} is a LIST not a dict
    \item Adjust batch size if GPU OOM errors occur
    \item Verify YOLOv8n model downloads successfully
\end{enumerate}

\section{Performance Optimization}

\subsection{GPU Memory Issues}

If encountering CUDA out of memory errors:
\begin{itemize}
    \item Reduce batch size to 8 or 4
    \item Reduce image size to 416 or 320
    \item Process frames at lower frequency (increase DETECTION\_INTERVAL)
    \item Disable unnecessary detection stages
\end{itemize}

\subsection{Slow Inference}

If detection is too slow:
\begin{itemize}
    \item Use YOLOv8n (nano) instead of larger models
    \item Enable detection interval (process every N frames)
    \item Reduce MJPEG stream quality
    \item Limit number of active zones
\end{itemize}

\section{Log Locations}

\begin{tabular}{|l|l|}
    \hline
    \textbf{Component} & \textbf{Log Location} \\
    \hline
    Backend FastAPI & stdout/stderr \\
    \hline
    Frontend Next.js & Browser console \\
    \hline
    Model Training & \texttt{training/runs/stage3\_fire\_smoke/} \\
    \hline
    Snapshots & \texttt{backend/static/snapshots/} \\
    \hline
    Video Jobs & Database video\_jobs table \\
    \hline
\end{tabular}

% ============================================================================
% APPENDIX
% ============================================================================
\appendix

\chapter{Glossary}

\begin{tabular}{|l|p{10cm}|}
    \hline
    \textbf{Term} & \textbf{Definition} \\
    \hline
    SIEWS & Smart Integrated Early Warning System \\
    \hline
    PPE & Personal Protective Equipment (helmet, vest, belt) \\
    \hline
    MJPEG & Motion JPEG - video streaming format \\
    \hline
    IoU & Intersection over Union - bounding box overlap metric \\
    \hline
    YOLO & You Only Look Once - object detection architecture \\
    \hline
    OCR & Optical Character Recognition - text reading from images \\
    \hline
    WebSocket & Bidirectional real-time communication protocol \\
    \hline
    ROI & Region of Interest \\
    \hline
    FPS & Frames Per Second \\
    \hline
    UUID & Universal Unique Identifier \\
    \hline
    Fonnte & WhatsApp API service provider \\
    \hline
\end{tabular}

\chapter{Class Enumerations}

\section{PPEClass}

\begin{lstlisting}[style=pythonstyle]
class PPEClass(IntEnum):
    UNKNOWN = 0
    BELT = 1
    HELMET = 2
    VEST = 3
\end{lstlisting}

\section{EnvClass}

\begin{lstlisting}[style=pythonstyle]
class EnvClass(IntEnum):
    BARRICADE = 0
    HARD_HAT = 1
    SAFETY_CONE = 2
    OPEN_HOLE = 3
    VEST_ENV = 4
\end{lstlisting}

\section{RoadClass}

\begin{lstlisting}[style=pythonstyle]
class RoadClass(IntEnum):
    LUBANG = 0    # Pothole
    RETAK = 1     # Crack
    TAMBALAN = 2  # Patch
\end{lstlisting}

\chapter{File Inventory}

\section{Backend Files}

\begin{longtable}{|l|l|p{6cm}|}
    \hline
    \textbf{File} & \textbf{Lines} & \textbf{Purpose} \\
    \hline
    main.py & ~800 & FastAPI app, all endpoints \\
    \hline
    stream.py & ~400 & StreamManager class \\
    \hline
    detector.py & ~300 & MultiStagePipeline \\
    \hline
    polygon.py & ~100 & Ray-casting algorithm \\
    \hline
    zone\_tracker.py & ~250 & PersonZoneTracker \\
    \hline
    violation\_checker.py & ~200 & ViolationChecker \\
    \hline
    notifier.py & ~100 & WhatsApp integration \\
    \hline
    shutdown.py & ~50 & Relay trigger \\
    \hline
    models.py & ~200 & SQLAlchemy models \\
    \hline
    database.py & ~50 & DB connection \\
    \hline
    face\_manager.py & ~300 & Face recognition \\
    \hline
    ocr\_engine.py & ~150 & OCR processing \\
    \hline
    detection\_models.py & ~100 & Class enums \\
    \hline
\end{longtable}

\section{Frontend Pages}

\begin{longtable}{|l|l|p{6cm}|}
    \hline
    \textbf{Page} & \textbf{Location} & \textbf{Purpose} \\
    \hline
    Dashboard & app/dashboard/page.tsx & Main monitoring interface \\
    \hline
    Incidents & app/incidents/page.tsx & Alert history and export \\
    \hline
    Zones & app/zones/page.tsx & Zone management \\
    \hline
    Settings & app/settings/page.tsx & System configuration \\
    \hline
    Analyze & app/analyze/page.tsx & Image analysis tool \\
    \hline
    Video Analysis & app/video-analysis/page.tsx & Video processing \\
    \hline
    Faces & app/faces/page.tsx & Personnel registry \\
    \hline
\end{longtable}

% References
\begin{thebibliography}{99}

\addcontentsline{toc}{section}{References}

bibitem{yolov8}
Ultralytics, ``YOLOv8 Documentation,'' 2024. [Online]. Available: \url{https://docs.ultralytics.com/}

bibitem{fastapi}
FastAPI, ``FastAPI Documentation,'' 2024. [Online]. Available: \url{https://fastapi.tiangolo.com/}

bibitem{nextjs}
Vercel, ``Next.js Documentation,'' 2024. [Online]. Available: \url{https://nextjs.org/docs}

bibitem{roboflow}
Roboflow, ``Roboflow Universe,'' 2024. [Online]. Available: \url{https://universe.roboflow.com/}

bibitem{fonnte}
Fonnte, ``Fonnte WhatsApp API,'' 2024. [Online]. Available: \url{https://fonnte.com/}

bibitem{dlib}
Davis King, ``dlib Documentation,'' 2024. [Online]. Available: \url{http://dlib.net/}

bibitem{easyocr}
JaidedAI, ``EasyOCR Documentation,'' 2024. [Online]. Available: \url{https://www.jaided.ai/easyocr/}

\end{thebibliography}

\printindex

\end{document}